\section{Experimental method: retain the actual measured support}
\label{sec:experimental-method}
\subsection{Datasets, signals and software responsibilities}
Two supplied MATLAB v7.3/HDF5 files are used. The outlier file contains
three rotor-reference slices at 2000 rpm, with $p=3$ and a constant used
resistance of \SI{7.695}{\milli\ohm}. The multi-speed file contains 1000 and
3000 rpm at a rotor reference of \SI{70}{\celsius}, with $p=4$ and a used
resistance of \SI{10}{\milli\ohm}. Thus the electrical speeds are
\SI{628.319}{\per\second} in the first experiment and
\SI{418.879}{\per\second}, \SI{1256.637}{\per\second} in the second.
These constants are read from each file, not inferred from the residuals.
Different pole-pair constants and unavailable machine provenance prohibit
pooling the two experiments into one resistance calibration.

The four supplied scripts define the method. The importer reads HDF5
references and exposes signals by name; it contains no machine physics.
\texttt{flux\_correction.py} implements voltage inversion, pair formation,
residuals and resistance-equivalent summaries. The outlier analyzer varies
rotor reference at one speed; the multi-speed analyzer varies speed at one
rotor reference and merges common requested keys. Both original analyzers
were executed before manuscript changes. Their calculation was preserved;
only headless/export options and a separate numeric/figure exporter were
added. No physics bug requiring a method correction was found.

The measured current signals are \texttt{PE\_Id\_meas\_avg} and
\texttt{PE\_Iq\_meas\_avg}. The voltage signals are
\texttt{PE\_Ud\_meas\_avg} and \texttt{PE\_Uq\_meas\_avg}.
The resistance and pole-pair constants are the
slice modes of \texttt{Const\_Machine\_R\_s} and
\texttt{Const\_Machine\_pole\_pairs}. Requested speed
\texttt{PE\_Rpm\_req} and \texttt{Rotor\_temp\_ref} select the block;
the requested speed also supplies \cref{eq:rpm-conversion}.
The supplied signal names establish access, but do not establish terminal
voltage calibration, averaging over whole mechanical revolutions, or the
upstream steady-state acceptance rule. These acquisition details remain to
be documented before an accuracy claim.

\subsection{Reduction, requested keys and measured denominators}
Within each requested-speed/rotor-reference slice, all samples sharing
\texttt{op\_index} are averaged. This index is supplied by the recorded
operating-point signal. The requested currents
\texttt{PE\_Id\_req\_meas} and \texttt{PE\_Iq\_req\_meas} are then rounded
to two decimal places solely to form geometric keys. The key spacing is
\SI{0.01}{\ampere}; each component moves by at most approximately
\SI{0.005}{\ampere}. It is a quantization rule, not a continuous distance
tolerance: nearly equal values straddling a rounding boundary can fail to
match. Measured currents and voltages are never rounded.

Repeated requested keys are aggregated by arithmetic means of their
operating-point rows. Consequently the second mean weights each
\texttt{op\_index} group equally, rather than weighting by its original sample
count. Flux inversion is linear at fixed slice resistance and speed, so
averaging flux before this key collapse is equivalent to inverting those
averaged signals with the same weights.
Positive and negative q keys are inner-joined at equal d key and opposite
q key. Zero-q keys and unmatched keys do not create pairs. Every retained
pair is an existing mirror pair \emph{in the requested-key geometry}; the
measured currents need not be mathematically mirrored. No missing partner
is interpolated, and no mirrored extrapolation supplies additional evidence.

Residuals use the paired reconstructed values, and
\cref{eq:measured-denominator} uses the actual measured pair magnitude $a_q$.
Map figures show flux at measured currents; pair figures use the requested
d key and measured $a_q$, specified in their captions.
Measured pair mismatch can leak allowed magnetic variation into the forbidden
channels. In the outlier slices the largest d-current mismatch is
\SIrange{4.024}{4.956}{\ampere}, and the largest opposite-q sum is
\SIrange{2.070}{2.249}{\ampere}. In the multi-speed slices these maxima are
\SI{0.606}{\ampere} and \SI{0.363}{\ampere} at 1000 rpm, and
\SI{0.506}{\ampere} and \SI{0.663}{\ampere} at 3000 rpm.
These deviations are recorded, not repaired using an unverified magnetic
reference. A future local mirror interpolation, for example with Delaunay
support, would require a separate bias/uncertainty study and is not implemented.

\subsection{Conditioning and descriptive summaries}
All exact-key pairs contribute to $r_d$ and $r_q$ summaries. The equivalent
residual is evaluated only for $a_q>0.10a_{q,\max}$, using the maximum over
pairs in the slice. This exclusion primarily prevents numerical amplification
of residual errors when dividing by small q-axis currents. It is not a
physical transition or a criterion proving that inverter effects vanish.

For an initial descriptive high-current summary, select
$0.85\leq a_q/a_{q,\max}\leq1$. The threshold is applied to measured
magnitudes, so membership may differ between slices even at identical keys.
Define the unscaled median absolute deviation by
\begin{equation}
 \boxed{\operatorname{MAD}(x)=
 \operatorname{median}\bigl(|x-\operatorname{median}(x)|\bigr).}
 \label{eq:raw-mad}
\end{equation}
The median and raw MAD limit the influence of isolated extreme values while
retaining the full plotted structure of the deliberately outlier-containing
dataset. No normal-distribution consistency multiplier is applied. MAD is
a descriptive spatial spread, not an uncertainty of the median or a
confidence interval; aggregated, correlated measurements are not independent
replicate calibrations. The high-current region is an empirical observation
of the present dataset, not a universal resistance-identification domain.

\subsection{Reproduction}
From the repository root, with Python, NumPy, pandas, SciPy, h5py and
Matplotlib installed, run the following. The commands preserve the default
pairing and current thresholds; omitting \texttt{--no-plots} opens the supplied
interactive plots, and \texttt{--export} generates manuscript data and
concrete vector-figure sources.
\begin{verbatim}
cd papers/flux_map_error_diagnostics/python
python analyze_outlier_dataset.py --no-plots --export
python analyze_multi_rpm_dataset.py --no-plots --export
\end{verbatim}
The exported CSV files retain full flux/pair tables, DAT files feed PGFPlots,
and JSON reports contain counts, package versions, settings, input and script
SHA-256 hashes, and descriptive metrics. Figure sources use only canonical
shared styles. No screenshot, fitted surface, smoothing, or outlier rejection
is required to regenerate the paper figures.
